\documentclass{situationist-preprint}

\title{Better to Have Been}
\msid{EXI-12-a}
\msfamily{EXI-12}
\msdiscipline{Meaning, Death and the Absurd}
\mstype{Principal manuscript}
\mskeywords{anti-natalism; Benatar; coming into existence; non-comparative benefit; quality of life; pessimism}

\begin{document}
\maketitle

\begin{abstract}
Is it better never to have been born? The pessimist tradition from Sophocles to Schopenhauer says yes, and Benatar has given the claim a rigorous form. ETH-20 shows that his asymmetry argument fails: evaluated on a single standard, never existing is better only for lives containing no good, and a life with more good than bad is better than never existing. This manuscript completes the answer. We argue, first, that the comparison can be made good in a sense that matters to the one who lives the life: although never existing is not a condition anyone is in, a life with more good than bad is good for the person who lives it, and to have been given it is a non-comparative benefit. Second, we answer Benatar's independent argument that our lives are much worse than we judge them to be. That argument contests which lives have more good than bad, not whether such lives are better than non-existence, and its case rests on an accounting that counts as bad what is merely not good. We conclude that a life containing more good than bad is better than non-existence.
\end{abstract}

\section{The problem}

The chorus of \emph{Oedipus at Colonus} says that not to be born is best, and next best to return quickly whence one came. Schopenhauer held that life is a business whose proceeds do not cover its costs. \citet{Benatar2006} argued for the stronger thesis that coming into existence is always a serious harm, by two independent routes. The first is the asymmetry: the absence of pain is good even if there is no one to enjoy it, while the absence of pleasure is not bad unless there is someone deprived of it. The second is the quality-of-life argument: well-known psychological biases, of optimism, adaptation and comparison, lead us to overrate the quality of our lives, which are in fact far worse than we judge.

ETH-20 addresses the first route. It shows that Benatar's asymmetry evaluates the absence of pain impersonally and the absence of pleasure personally, that on any single standard never existing comes out better only for lives without any good, and that parity requires the impersonal standard, on which the absence of pleasure, though bad for no one, is not neutral. We take this result as established and do not repeat it. Two questions remain for the conclusion of the present family. The first is in what sense a life is \enquote{better than non-existence}, since no one is in the condition of never existing. The second is whether Benatar's second route reinstates his conclusion.

\section{Framework}

\begin{definition}[Balance]\label{d:bal}
A life has \emph{more good than bad} if the goods it contains for its subject (pleasures, satisfied concerns, achievements, relationships and other components of well-being) outweigh the bads it contains for her (pains, frustrations, losses).
\end{definition}

\begin{principle}[Impersonal comparison]\label{p:imp}
Of two outcomes alike except that in one a person exists, the one in which she exists is impersonally better if her life has more good than bad (ETH-20).
\end{principle}

\begin{principle}[Non-comparative benefit]\label{p:noncomp}
An event benefits a person non-comparatively if it causes her to have a life that is good for her. A non-comparative benefit does not require that she would have been worse off without it \citep{Holtug2001, Harman2009}.
\end{principle}

Principle~\ref{p:noncomp} answers the familiar worry that coming into existence cannot benefit anyone because there is no prior state in which the beneficiary was worse off. The worry relies on a comparative notion of benefit. But we commonly judge that a person's life is good for her, that a pleasure is good for the one who has it, without asking what her condition would have been without it, and these judgments do not lapse because the alternative is one in which she does not exist. The deprivation account of the badness of death, defended in EXI-02, already accepts that a subject's good can be assessed against an alternative in which she is not there to be worse off \citep{Nagel1970, Feldman1991}.

\section{The argument}

\begin{theorem}\label{t:main}
A life containing more good than bad is better than non-existence: impersonally, and in that it is a non-comparative benefit to the one who lives it.
\end{theorem}

\begin{proof}
The impersonal comparison is Principle~\ref{p:imp}, established in ETH-20. For the second clause: a life with more good than bad is good for its subject by Definition~\ref{d:bal}, since the goods it contains for her outweigh the bads. By Principle~\ref{p:noncomp}, to be caused to have such a life is a non-comparative benefit. So the life is better than non-existence in both senses that the question can bear: the world in which it is lived is better, and the person who lives it has received something good for her.
\end{proof}

\begin{proposition}\label{p:pollyanna}
Benatar's quality-of-life argument does not reinstate the thesis that never existing is better for lives with more good than bad.
\end{proposition}

\begin{proof}
The argument claims that our judgments of our lives' quality are inflated, and that lives are worse than we think. If sound, it lowers our estimate of how many lives satisfy the condition of Definition~\ref{d:bal}. It does not bear on the conditional established by Theorem~\ref{t:main}, which concerns lives that do satisfy it. To yield the conclusion that \emph{every} life is worse than non-existence, the argument would need to show that no life has more good than bad. Benatar's case for that stronger claim rests on counting among a life's bads the ever-present gaps between desire and satisfaction, the bare absence of more pleasure than one has. These are not bads in the life but goods it does not contain; to count them as bads is to apply to the existing person the very evaluation of absences that ETH-20 shows to be incompatible with Benatar's own evaluation of absent pains. Without that accounting, the psychological biases he cites show only that estimates should be corrected, not that the corrected balance is always negative.
\end{proof}

\section{Objections and replies}

\begin{objection}[Bad for no one]\label{o:noone}
If I had never existed, there would have been no one for whom that was bad. So existence cannot be better \emph{for me} than non-existence.
\end{objection}

\begin{reply}
Theorem~\ref{t:main} does not say that existence is comparatively better for me than non-existence would have been. It says that my life is good for me, which is a non-comparative claim about me, and that the world containing it is impersonally better. Both claims are true even though non-existence would not have been bad for anyone.
\end{reply}

\begin{objection}[Testimony of the suffering]\label{o:test}
Some people sincerely judge that it would have been better had they never been born. Their testimony should count.
\end{objection}

\begin{reply}
It counts, and for those whose lives contain more bad than good, Theorem~\ref{t:main} does not apply to them; nothing in our argument denies that some lives are worse than non-existence. The thesis is conditional on the balance of the life. What the testimony does not show is that lives with more good than bad are worse than non-existence, which is what the anti-natalist must establish.
\end{reply}

\begin{objection}[Risk]\label{o:risk}
Even if a good life is better than never existing, no one can guarantee that a life will be good. Creating a person imposes a risk on someone who cannot consent \citep{Benatar2006}.
\end{objection}

\begin{reply}
The objection concerns the ethics of procreation under uncertainty, not the comparison of a life with non-existence. It is answered in the conclusion of ETH-20, that having children whose lives will be good on balance is not wrong. The question of this family is settled by Theorem~\ref{t:main}, whatever the right policy for imposing risks on future persons.
\end{reply}

\section{Conclusion}

It is not better never to have been born, for a life containing more good than bad. Such a life makes the world better than it would have been without it, as ETH-20 shows, and it is good for the one who lives it, which is a benefit to her though there was no one who would otherwise have been worse off. Benatar's argument from the poor quality of life contests which lives have more good than bad, and it does so by counting as bads what are only the limits of a life's goods. Some lives are worse than non-existence; lives with more good than bad are better.

\begin{thebibliography}{99}\small
\bibitem[Benatar(2006)]{Benatar2006} Benatar, D. (2006). \emph{Better Never to Have Been: The Harm of Coming into Existence}. Clarendon Press.
\bibitem[Feldman(1991)]{Feldman1991} Feldman, F. (1991). Some puzzles about the evil of death. \emph{The Philosophical Review}, 100(2), 205--227.
\bibitem[Harman(2009)]{Harman2009} Harman, E. (2009). Critical study: David Benatar, \emph{Better Never to Have Been}. \emph{No\^us}, 43(4), 776--785.
\bibitem[Holtug(2001)]{Holtug2001} Holtug, N. (2001). On the value of coming into existence. \emph{The Journal of Ethics}, 5(4), 361--384.
\bibitem[Nagel(1970)]{Nagel1970} Nagel, T. (1970). Death. \emph{No\^us}, 4(1), 73--80.
\end{thebibliography}

\end{document}
